\section{Native compilation, bit entries and the quantitative bound}
\label{sec:bridges}

This section assembles \cref{thm:cascades}. The native compiler and the
canonical-width C5 assembly are formalized; the bit-to-word bridge below
has a written proof.
The bit-to-word construction is the parallel oblivious simulation already
described by Harvey and van der Hoeven \cite[Section~7.6]{HarveyHoeven2025}.
We check its motion and correctness-set accounting because a time bound
alone would not transfer the affine-motion obstruction.

\subsection{Time-dependent gates cost only logarithmically many clock bits}

The exact compiler is formalized independently of C5.
\formal{CascadeObstruction/CompilerCore.lean}{CascadeObstruction.Compiler.run_denote,
CascadeObstruction.Compiler.stepN_eq,CascadeObstruction.Compiler.haltsAt_iff,
CascadeObstruction.Compiler.state_count,CascadeObstruction.Compiler.state_bits_le}
\begin{outline}
For a completion of duration $T$, form a native finite controller with
states $(t,s)$, $0\le t\le T$, $s\in Q$. At $t<T$, use the prescribed
local rule, movements and append event, then increment $t$. At $t=T$,
halt. This has at most $(T+1)|Q|$ states and hence at most
\[
 q+\lceil\log_2(T+1)\rceil
\]
controller bits. Each source tick is one actual native transition; the
old tuple is read before any of its work symbols is overwritten. The
input, work tapes, output and their positions agree step by step. There
is no oracle for time and no free access. The native transition table
may depend on $d$, exactly as allowed by Paper~2's nonuniform theorem.
Since $T=O(L)$, the added bit cost is $O(\log L)=o(L)$.
\end{outline}
\begin{samepage}
Prescribed head moves and append times are preserved on every input,
with arbitrary simultaneous local data updates.
\formal{CascadeObstruction/CompilerSchedule.lean}{CascadeObstruction.Compiler.scheduled_work_trajectory,
CascadeObstruction.Compiler.scheduled_input_trajectory,CascadeObstruction.Compiler.scheduled_append,
CascadeObstruction.Compiler.scheduled_oblivious}
\end{samepage}
\begin{outline}
Induct on the public tick index. Each input and work position changes by
the prescribed unit move, independently of the data rule. The append
decision is the prescribed Boolean. After tick $T$, the native machine
halts and its trajectories freeze.
\end{outline}

At canonical width, suppose along some $d_i\to\infty$ both
$q_i/L_i\to0$ and $\deficit_i/L_i\to0$. Compile these completions, set
$e_i=\lceil\deficit_i\rceil$, and apply \cref{thm:imported} to the
all-regular clocks of \cref{sec:schedules}. Its conclusion would be
$1/(2G)\le0$, a contradiction. Padding a smaller alphabet into one
fixed alphabet, if necessary, fixes its cell charge $G$.

\subsection{An oblivious bit-to-word lift preserving the limiting motion}

Let $N=d^2$. Consider a bit-entry completion taking $T=O(N)$ ticks,
with $k$ fixed work tapes, alphabet $\Gamma$, controller $Q$, and
correctness set $S\subseteq\{0,1\}^N$. Put $b=\log_2d$ and $L=bN$.
The canonical word input already consists of interleaved bit planes:
coordinate $p$ of plane $\ell$ is at raw address $bp+\ell$.
Replace every old work cell $p$ by its block of $b$ cells
$bp,\ldots,bp+b-1$, on the same work tape over the same alphabet.
Keep $b$ old controllers in the new finite control, plus a public clock.

An old tick is simulated in exactly $3b$ native ticks. Initially every
head is at the left endpoint $bp$ of its old cell's block. In the first
$b$ ticks, visit lanes $\ell=0,\ldots,b-1$ from left to right. Read lane
$\ell$'s controller, its raw bit and all its old work symbols, apply the
old local rule, and update that lane only. Emit its bit if this old tick
is an append event. Advance all heads after the first $b-1$ lane ticks
and stay after the last. Every old read tuple is obtained before that
lane's simultaneous writes.

Every head is now at $bp+b-1$. If its old prescribed move was
$\Delta\in\{-1,0,1\}$, its next target is $b(p+\Delta)$, at distance
at most $2b-1$. Use the next $2b$ ticks to reach it, parking when it
arrives. These transit ticks preserve every scanned work symbol and
all lane controllers and emit nothing. Thus all heads reach the next
block's left endpoint together. Intermediate reads grant no additional
data access: they are ignored except for preserving the symbol read.
Negative and right-exterior old raw cells expand to entirely exterior
blocks. Induction proves exact simulation of every lane.

The resulting native program has unchanged alphabet and tape count,
time $3bT$, and at most
\[
 (3bT+1)|Q|^b\quad\text{states},\qquad
 bq+\lceil\log_2(3bT+1)\rceil\quad\text{controller bits}.
\]
Each old output coordinate emits all $b$ lanes consecutively in their
correct order. The lift is correct on the interleavings of the product
set $S^b$; there are $|S|^b$ of them.
Their deficit is $b(N-\log_2|S|)$. No preliminary bit-plane rearrangement
is used.

Let $h$ be one original trajectory, linearly interpolated between ticks,
and $h'$ its lifted trajectory. Within the block for old tick $t$,
$h'$ stays within $2b$ of $bh(t)$. The original interpolation changes
by at most one cell over that tick. Thus, with frozen continuations,
\begin{equation}
\label{eq:lift-error}
 \left|\frac{h'(\tau)}{bN}
        -\frac{h(\tau/(3b))}{N}\right|\le\frac3N
 \qquad(\tau\ge0).
\end{equation}
If the original physical paths converge to $f(s)$, the lifted physical
paths converge to $f(s/3)$. A locally affine limit outside a null set
retains that property under this dilation.

The geometric clock proof in \cref{sec:schedules} holds at $b=1$.
For any finite retained depth, its stage endpoints and profiles converge.
The deep middle interval has duration tending to zero uniformly as the
retained depth increases, by \eqref{eq:tail}. Unit-Lipschitz continuity
controls the displacement it can contribute. Consequently the complete
original paths have locally affine limits outside their countable piece
endpoints and the recursive accumulation time. The frozen halt tail is
constant. Equation~\eqref{eq:lift-error} gives the same property to the
lifted limits.
If the original bit completions had $q+\deficit=o(N)$, their lifts
would satisfy the canonical resource and near-full correctness
assumptions. Paper~2's native affine-time bound
\eqref{eq:affine-toll} is then contradicted. The lifted trace need
\emph{not} equal the canonical-width cascade constructor. It only needs
the already established limiting-motion property. Oblivious head and
append schedules are essential to this simulation.

\subsection{A positive uniform deficit}

\begin{lemma}[Sequential uniform-gap criterion]
\label{lem:uniform-gap}
For each natural size $n$, let $\mathcal A_n$ be any type of admissible
objects and let $c_n:\mathcal A_n\to[0,\infty)$ be a cost. Suppose that
for every sequence $n_i\to\infty$ and $x_i\in\mathcal A_{n_i}$, the costs
$c_{n_i}(x_i)$ do not converge to zero. Then there are $\delta>0$ and $N$
such that $c_n(x)\ge\delta$ for all $n\ge N$ and all $x\in\mathcal A_n$.
\end{lemma}
\formal{CascadeObstruction/UniformGap.lean}{CascadeObstruction.exists_uniform_positive_gap}
\begin{outline}
Failure supplies $n_i\ge i$ and $x_i$ with
$0\le c_{n_i}(x_i)<1/(i+1)$. The sizes diverge and the costs tend to zero
by squeezing. No compactness of the object types is required.
\end{outline}

\begin{corollary}[Native regular-clock classes]
\label{cor:native-gap}
Consider a size-indexed class of native machines with canonical input
length $L(n)$, fixed work-tape count and cell charge, and fixed $a\ge0$.
Each instance supplies a correct-input set $W$, an integer deficit $e$
with $2^{L(n)}\le2^e|W|$, and a complete unary two-stage clock with
parameters $(\rho,C)=(1/2,1)$ for each of its correct halting executions.
Suppose every diverging sequence chosen from the class has the full
family-level regularity of \cref{sec:criterion}. Then, for one $\delta>0$
and all sufficiently large $n$, every instance in the class satisfies
\[
 \frac{\lceil\log_2 Q_{\rm native}\rceil+e}{L(n)}\ge\delta,
\]
where $Q_{\rm native}$ is its number of controller states.
\end{corollary}
\formal{CascadeObstruction/NativeObstruction.lean}{CascadeObstruction.no_vanishing_native_dyadic_cost,CascadeObstruction.uniform_native_dyadic_gap}
\begin{outline}
Vanishing combined cost forces both nonnegative summands to vanish.
Paper~2's \lean{no_native_regular_recursion} rules out that sequence.
Apply \cref{lem:uniform-gap}. In Lean, the machine has
$\texttt{TM.q}+1$ states and the bit count is
$\texttt{Nat.clog 2 (TM.q+1)}$. The integer deficit and clock witnesses
are fields of \lean{CascadeObstruction.NativeDyadicInstance}; regularity
remains an explicit hypothesis on the class. ``Dyadic'' here describes
the clock contraction, not a restriction on the native matrix dimension.
\end{outline}

\begin{informalproof}
To finish the proof of \cref{thm:cascades}, fix an alphabet, one constructor
and one of the two widths. The preceding arguments rule out every
diverging sequence of completions with $(q+\deficit)/L\to0$.
If no positive uniform gap held eventually, one could choose strictly
increasing legal dyadic dimensions $d_j$ and completions with
\[
                       0\le\frac{q_j+\deficit_j}{L_j}<\frac1{j+1}.
\]
Both nonnegative summands would tend to zero. Rounding the deficit and
adding the public clock preserve this convergence, giving the forbidden
sequence. Hence some positive gap and finite threshold exist for this
constructor and width. Take the minimum of the ten positive gaps and
the maximum of their thresholds. This proves \cref{eq:main} simultaneously.
The sequential-to-uniform step is also checked in the new Lean library;
the remaining schedule and simulation premises are supplied by the written
proofs. For C5 at canonical width, all these premises are also discharged
in Lean below.
\end{informalproof}

This argument does not estimate the gap or threshold. Finite experiments
cannot turn its existential constants into certified numerical ones.

\subsection{The complete canonical-width C5 formalization}
\label{sec:c5-formal}

The C5 part of \cref{thm:cascades}, at canonical width, is checked as a
theorem about a specified finite tick word, not about a supplied
regularity certificate.
\formal{CascadeObstruction/C5Exclusion.lean}{CascadeObstruction.C5.uniform_c5_gap,
CascadeObstruction.C5.completion_family_regular}
\begin{outline}
The type \lean{CascadeObstruction.C5.CorrectCompletion} contains a legal
dimension $d=2^{K+1}$ with $K\ge1$, a finite controller, arbitrary local
data rules attached to the C5 tick word, and a nonempty set of correctly
transposed inputs. It contains no clock, approximation or regularity
hypothesis.

The tick word concatenates root preparation, forward stages $1,\ldots,K$,
backward stages $K,\ldots,1$, and root output. Distinct stage labels and
prefix sums construct a complete partition of the actual tick indices.
The exact bounds in \cref{sec:schedules} give the clock with $a=18$.
Extracting each contiguous stage from the full word identifies its
frozen displacement with the block sequence used in the ramp argument.
The resulting regularity theorem applies to all positive widths and all
eligible subsequences of actual stages.

The compiler then supplies a native machine with the same work alphabet,
two work tapes, the same symbols and output, and precisely the prescribed
head paths. Its additional controller cost is at most
$\lceil\log_2(36L+1)\rceil$, uniformly over the local rules.
For $D=L-\log_2|S|$, rounding to $e=\lceil D\rceil$ gives
$2^L\le2^e|S|$ and $e<D+1$. Apply the native uniform-gap theorem and
absorb the clock cost and this one-bit rounding error into a smaller
positive gap. This proves the real-deficit conclusion with all model
and schedule obligations discharged.
\end{outline}

The reusable clock, regularity and deficit-transfer interfaces are recorded
by the following declarations.
\formal{CascadeObstruction/C5Clock.lean}{CascadeObstruction.C5.clock,
CascadeObstruction.C5.trace_time_bound}
\formal{CascadeObstruction/C5Regularity.lean}{CascadeObstruction.C5.clock_regular}
\formal{CascadeObstruction/RasterRegularity.lean}{CascadeObstruction.Raster.curve_subsequence,
CascadeObstruction.Raster.regular_of_ramp_approximation}
\formal{CascadeObstruction/GapTransfer.lean}{CascadeObstruction.uniform_native_real_deficit_gap}
\begin{outline}
The clock declaration packages the proved partition and budgets. The
regularity proof combines finite-product compactness for ramp parameters
with a vanishing-error estimate. The deficit-transfer theorem rounds
the correctness deficit internally and removes any uniformly sublinear
compiler overhead. These interfaces reuse Paper~2's obstruction.
\end{outline}

\subsection{Public initial words and paid entry information}

\begin{proposition}[Complete entry information]
\label{prop:entry}
The same five-schedule bound holds, with possibly smaller $\delta>0$,
if arbitrary public initial work words are allowed. More generally, let
$Z(X)$ specify the entire initial work contents at all visited cells and
the initial controller, with the same zero starting head positions, empty output and one
common set of later rules. Under the uniform law on a nonempty correctness
set $S$,
\begin{equation}
\label{eq:entry}
                 q+L-\log_2|S|+H(Z)\ge\delta L.
\end{equation}
\end{proposition}
\notformalized
\begin{informalproof}
The complete schedule visits only an $O(L)$ coordinate range on each
tape. A constant number of outward and return scans installs any public
initial word and restores all starting heads. Hardwire the public symbols
in the nonuniform transition table, ignore the raw input during these
scans, and install the prescribed initial controller. The time is $O(L)$,
the additional clock has $O(\log L)$ bits, and the extra profiles are
finitely affine. The preceding obstruction and sequential argument apply
uniformly over all public words. The bit case uses the same lane lift.

For input-dependent entry information, fix each nonempty fibre $S_z$.
Its complete entry is now public, so
$\log_2|S_z|\le(1-\delta)L+q$ with the same constants. Averaging under
uniform $X\in S$ gives $H(X\mid Z)\le(1-\delta)L+q$.
As $Z$ is a deterministic function of $X$,
$\log_2|S|=H(X)=H(X\mid Z)+H(Z)$, proving \eqref{eq:entry}.
\end{informalproof}

Thus public setup or sublinear joint entry information cannot rescue the
same core. A precomputed transpose supplies $L$ bits of entry information
and is not excluded by this statement. It still needs a paid producer.
An output prefix, input-dependent starting heads, or an input-dependent
choice of later rules is not part of this interface; none may be omitted
from the information bill by calling it preparation.
